\section{Historia zmian}\label{historia-zmian}

\begin{longtable}[]{@{}
  >{\raggedright\arraybackslash}p{(\columnwidth - 4\tabcolsep) * \real{0.2}}
  >{\raggedright\arraybackslash}p{(\columnwidth - 4\tabcolsep) * \real{0.2}}
  >{\raggedright\arraybackslash}p{(\columnwidth - 4\tabcolsep) * \real{0.6}}@{}}
\toprule\noalign{}
\begin{minipage}[b]{\linewidth}\raggedright
Wersja
\end{minipage} & \begin{minipage}[b]{\linewidth}\raggedright
Data
\end{minipage} & \begin{minipage}[b]{\linewidth}\raggedright
Zmiany
\end{minipage} \\
\midrule\noalign{}
\endhead
\bottomrule\noalign{}
\endlastfoot
1.2 & 2026-04-03 & \textbf{Nowe hasła --- bojowe wozy desantowe (nowa
kategoria):} BMD-1, BMD-2, BMD-3, BMD-4. \textbf{Nowe hasła ---
transportery opancerzone:} MT-LB. Łącznie dodano 5 haseł (130→135),
liczba kategorii wzrosła z 16 do 17. \\
1.1 & 2026-04-03 & \textbf{Nowe hasła --- czołgi:} T-55, T-62, T-62M,
2S25 Sprut-SD. \textbf{Nowe hasła --- artyleria samobieżna:} 2S19
Msta-S. \textbf{Nowe hasła --- systemy przeciwlotnicze:} ZU-23-2,
ZSU-23-4 Szyłka, S-60, 2K22 Tunguska, 9K33 Osa. Łącznie dodano 10 haseł
(120→130). \textbf{Poprawki:} zdjęcie T-62M (było identyczne jak T-62,
zastąpiono fotografią z targów Army-2023); zdjęcie miny VS-50 (błędnie
pokazywało brazylijską rakietę, zastąpiono fotografią miny). \\
1.0 & 2026-04-01 & Pierwsze wydanie. 120 haseł w 16 kategoriach: czołgi
(3), BWP (3), APC (3), artyleria samobieżna (7), artyleria holowana
(11), systemy przeciwlotnicze (7), ATGM (5), lotnictwo (5), karabiny
szturmowe (10), karabiny maszynowe (7), karabiny snajperskie (7), RPG i
wyrzutnie (11), wyrzutnie granatów (7), pistolety i PM (15), granaty
ręczne (6), miny (13). 134 obrazy pobrane z Wikimedia Commons. \\
\end{longtable}
